\chapter{进程调度}

\section{调度基本概念}
\concept{调度队列}
\begin{points}
  \item 作业队列：外存上等待进入系统的作业集合。
  \item 就绪队列：内存中已就绪、等待 CPU 的进程集合。
  \item 设备队列：等待某个 I/O 设备的进程集合。
  \item 进程在整个生命周期中会在这些队列之间迁移。
\end{points}
\plain{OS 把进程/作业放进不同队列，调度就是从队列里按规则挑下一个。}
\exam{常考作业队列/就绪队列/设备队列，高级/中级/低级调度分别管什么。}


\concept{三级调度}
\begin{points}
  \item 高级调度又称作业调度/长程调度，决定哪些作业从外存进入内存，影响多道程序度。
  \item 中级调度又称交换调度，决定哪些进程换入、换出内存，用于缓解内存紧张。
  \item 低级调度又称进程调度/短程调度，决定就绪队列中哪个进程获得 CPU，运行频率最高。
  \item 记忆：高级管“进不进系统”，中级管“在不在内存”，低级管“谁上 CPU”。
\end{points}
\plain{OS 把进程/作业放进不同队列，调度就是从队列里按规则挑下一个。}
\exam{常考作业队列/就绪队列/设备队列，高级/中级/低级调度分别管什么。}


\concept{CPU 密集型与 I/O 密集型}
\begin{points}
  \item CPU 密集型作业计算时间多，CPU burst 长，I/O 次数相对少。
  \item I/O 密集型作业 I/O 操作多，CPU burst 短。
  \item 作业调度应合理搭配两类作业：CPU 密集型过多会导致 I/O 设备空闲；I/O 密集型过多会导致就绪队列空、CPU 空闲。
\end{points}
\plain{硬件基础不用背太深，重点是知道哪些硬件机制支撑 OS：寄存器保存现场，PSW 保存状态，CPU 模式限制权限。}
\exam{常考 CPU 工作流程、单/多处理器、寄存器类别、PSW 中断位/状态位。}


\concept{调度程序与分派程序}
\begin{points}
  \item 调度程序负责从就绪队列中选择下一个运行进程。
  \item 分派程序负责完成上下文切换、切换到用户态、跳转到被选进程的正确位置。
  \item 分派延迟是停止一个进程并启动另一个进程所需时间，应尽可能短。
\end{points}
\plain{把这个知识点放回本章主线理解：它回答的是 OS 为了管理资源、隐藏硬件、保证并发正确性而采用的某个对象、规则或步骤。}
\exam{常考选择/填空/简答，让你判断概念归属、比较相近概念，或按“定义-作用-机制-易错点”展开。}


\section{调度目标与评价指标}
\concept{常见评价指标}
\begin{points}
  \item CPU 利用率：CPU 忙碌时间比例，越高越好。
  \item 吞吐量：单位时间完成的进程数量，越高越好。
  \item 周转时间：进程从提交到完成的总时间，越短越好。
  \item 等待时间：进程在就绪队列中等待 CPU 的总时间，越短越好。
  \item 响应时间：从提交请求到首次响应的时间，分时系统尤其关注。
\end{points}
\plain{调度题本质是按时间轴算完成时间，再由完成时间反推周转、等待、响应。}
\exam{常考平均周转时间、带权周转时间、等待时间计算，务必画甘特图。}


\concept{周转时间公式}
\begin{points}
  \item 完成时间 $C_i$，到达时间 $A_i$，服务时间 $S_i$。
  \item 周转时间 $T_i=C_i-A_i$。
  \item 带权周转时间 $W_i=\dfrac{T_i}{S_i}$。
  \item 平均周转时间和平均带权周转时间就是对所有进程取平均。
\end{points}
\plain{调度题本质是按时间轴算完成时间，再由完成时间反推周转、等待、响应。}
\exam{常考平均周转时间、带权周转时间、等待时间计算，务必画甘特图。}


\section{调度方式}
\concept{抢占式与非抢占式}
\begin{points}
  \item 非抢占式调度：进程一旦获得 CPU，除非主动放弃、阻塞或结束，否则不会被强行剥夺。
  \item 抢占式调度：操作系统可因时间片到、更高优先级进程到达等原因剥夺 CPU。
  \item 分时系统通常需要抢占式调度；批处理系统可以使用非抢占式调度。
\end{points}
\plain{抢占就是正在运行的进程也可能被更高优先级/时间片到期打断。}
\exam{常考抢占式与非抢占式适用场景。}


\section{常用调度算法}
\concept{先来先服务 FCFS}
\begin{points}
  \item 按进程到达就绪队列的先后顺序分配 CPU。
  \item 优点：简单、公平、不会饥饿。
  \item 缺点：短作业可能排在长作业后面，平均等待时间可能很长，存在“护航效应”。
  \item 适合批处理场景，不适合交互式短任务较多的场景。
\end{points}
\plain{谁先来谁先上，最朴素但容易让短作业排在长作业后面。}
\exam{常考护航效应和平均等待时间计算。}


\concept{短作业优先 SJF 与最短剩余时间优先 SRTF}
\begin{points}
  \item SJF 选择预计运行时间最短的作业先执行，是非抢占式算法。
  \item SRTF 是抢占式版本，总是选择剩余运行时间最短的进程。
  \item 优点：理论上可使平均等待时间最短。
  \item 缺点：需要预测运行时间；长作业可能长期等待，产生饥饿。
\end{points}
\plain{优先做短的，平均等待时间好看，但长作业可能一直等。}
\exam{常考非抢占 SJF 和抢占式 SRTF 的甘特图区别。}


\concept{优先级调度}
\begin{points}
  \item 为每个进程赋予优先级，优先级高者先运行。
  \item 可分为抢占式和非抢占式。
  \item 静态优先级创建后不变，动态优先级可根据等待时间、资源使用情况改变。
  \item 缺点是低优先级进程可能饥饿；常用老化 aging 提高等待过久进程的优先级。
\end{points}
\plain{给进程分等级，重要的先跑；问题是低优先级可能饥饿。}
\exam{常考静态/动态优先级、抢占/非抢占、老化技术防饥饿。}


\concept{时间片轮转 RR}
\begin{points}
  \item 每个就绪进程轮流获得一个时间片，时间片用完仍未完成则回到就绪队列尾部。
  \item 适合分时系统，强调响应时间和公平性。
  \item 时间片太大，退化为 FCFS；时间片太小，上下文切换开销过大。
  \item 做题时严格画时间轴，注意新到达进程插入队列的位置和时间片剩余规则。
\end{points}
\plain{每人轮流跑一小段，适合分时系统；时间片太大像 FCFS，太小切换开销大。}
\exam{常考时间片轮转甘特图，以及时间片大小对响应时间和开销的影响。}


\concept{高响应比优先 HRRN}
\begin{points}
  \item 响应比公式：$R=\dfrac{等待时间+服务时间}{服务时间}=1+\dfrac{等待时间}{服务时间}$。
  \item 等待时间越长，响应比越大；服务时间越短，响应比也越大。
  \item 该算法兼顾短作业优先和防止长作业饥饿。
  \item 通常是非抢占式调度，用在作业调度中较常见。
\end{points}
\plain{每人轮流跑一小段，适合分时系统；时间片太大像 FCFS，太小切换开销大。}
\exam{常考时间片轮转甘特图，以及时间片大小对响应时间和开销的影响。}


\concept{多级队列调度}
\begin{points}
  \item 将就绪队列划分为多个队列，不同队列服务不同类型进程，如系统进程、交互进程、批处理进程。
  \item 队列之间有固定优先级或时间比例分配，队列内部可使用不同调度算法。
  \item 固定多级队列的问题是低优先级队列可能长期得不到服务。
\end{points}
\plain{新来的先给高优先级小时间片；用得越久越往下掉，适合交互和批处理混合。}
\exam{常考多级反馈队列调度过程，注意抢占、降级、时间片没用完是否保留。}


\concept{多级反馈队列 MLFQ}
\begin{points}
  \item 多级反馈队列允许进程在不同优先级队列之间移动。
  \item 新进程通常先进入高优先级队列，时间片较短，若用完时间片还未完成则降级。
  \item 等待过久的低优先级进程可被提升，避免饥饿。
  \item 直觉：先假设新来的任务可能是短交互任务，给它快速响应；若它表现得像长计算任务，就逐步降到低优先级。
\end{points}
\plain{队列是 OS 组织进程的方式：就绪的排就绪队列，等设备的排设备队列。}
\exam{常考进程在生命周期中如何在不同队列迁移。}


\section{实时调度、多处理器调度与线程调度}
\concept{实时调度基本概念}
\begin{points}
  \item 实时系统中，正确结果如果迟到，也可能等同于错误结果。
  \item 周期性任务常用处理时间 $t_i$、截止期限 $d_i$、周期 $p_i$ 描述，通常 $0\le t_i\le d_i\le p_i$。
  \item 必要的可调度条件之一：$\sum_{i=1}^{m}\dfrac{t_i}{p_i}\le 1$，表示 CPU 总利用需求不能超过 100\%。
\end{points}
\plain{实时调度不只看快慢，而是看能否在截止期限前完成。}
\exam{常考硬/软实时、中断延迟/分派延迟、RMS/EDF 和可调度条件。}


\concept{速率单调 RMS 与最早截止期限 EDF}
\begin{points}
  \item RMS 是静态优先级算法，周期越短，优先级越高。
  \item EDF 是动态优先级算法，截止期限越近，优先级越高。
  \item EDF 在单处理器、理想条件下可达到更高利用率，但实现和开销更复杂。
\end{points}
\plain{实时调度不只看快慢，而是看能否在截止期限前完成。}
\exam{常考硬/软实时、中断延迟/分派延迟、RMS/EDF 和可调度条件。}


\concept{多处理器调度}
\begin{points}
  \item 非对称多处理调度中，一个主处理器负责调度，其他处理器执行任务。
  \item 对称多处理调度中，各处理器都可参与调度，常见挑战是负载均衡和缓存亲和性。
  \item 处理器亲和性指尽量让进程继续在原 CPU 上运行，以利用缓存中的数据。
\end{points}
\plain{硬件基础不用背太深，重点是知道哪些硬件机制支撑 OS：寄存器保存现场，PSW 保存状态，CPU 模式限制权限。}
\exam{常考 CPU 工作流程、单/多处理器、寄存器类别、PSW 中断位/状态位。}


\concept{线程调度}
\begin{points}
  \item 用户级线程由线程库调度，内核只看到进程。
  \item 内核级线程由内核调度，可以在多个处理器上并行。
  \item 线程调度更强调轻量切换和同一进程内线程之间的同步关系。
\end{points}
\plain{线程是进程里面更轻的执行流，同进程线程共享资源，但各自有栈、寄存器和 PC。}
\exam{常考进程/线程区别、用户级线程和内核级线程优缺点、多对一/一对一/多对多模型。}


